

\section{Conclusion}

\textbf{Conclusion.} 
Ling 2.0 demonstrates that large-scale sparse language foundations can advance both reasoning capability and computational efficiency through coordinated innovations in architecture, training, and infrastructure. With its high-sparsity Mixture-of-Experts design, reasoning-oriented data pipeline, multi-stage alignment strategy, and FP8-based trillion-scale infrastructure, Ling 2.0 establishes a scalable foundation for general reasoning models. The three released models—\textbf{\texttt{Ling-mini-2.0}}, \textbf{\texttt{Ling-flash-2.0}}, and \textbf{\texttt{Ling-1T}}—consistently follow the Ling Scaling Law and collectively define a new Pareto frontier between reasoning accuracy and computational cost, illustrating the effectiveness of the ``every activation boosts'' principle.  

Despite these advances, Ling 2.0 still faces several open challenges. 
First, its current grouped-query attention (GQA) architecture constrains efficiency in long-context scenarios; ongoing work explores linear and sparse-attention designs to further improve scalability. 
Second, while Ling 2.0 achieves strong reasoning precision and efficiency, the effective reasoning length and depth still have room for enhancement. 
Finally, complex instruction following and agentic behaviors remain under development. Building upon Ling 2.0's strong reasoning foundation, future work will extend toward more general, autonomous, and interactive capabilities.  

Together, these directions mark the next step in scaling general intelligence---toward models that not only think more efficiently, but also act more generally.

